%=============================================================================!
% 4.10  MOLECULAR VIBRATIONS AND THE TIME STEP                                !
%=============================================================================!
\section{Molecular vibrations and the time step}
\label{sec:os-molecules}

The oscillations of this chapter are the motions of atoms in molecules.
This section turns measured vibrational frequencies into periods in
femtoseconds and draws the consequence for simulation: a computer that
advances the motion in steps must take steps short compared with the
fastest of these periods.

\subsection*{Frequencies from light}

Light is an oscillation that travels, at the speed of light $c$; its
wavelength is the distance from one crest to the next, and infrared light
has wavelengths too long for the eye to see. A molecule absorbs infrared
light at frequencies close to those of many of its normal modes, much as
a driven oscillator responds most strongly at resonance
(\cref{sec:os-driven}), and other kinds of spectrum reveal most of the
rest, so measured spectra give vibrational frequencies. Spectroscopists
quote them as wavenumbers $\tilde\nu$, the number of wavelengths of the
light per centimetre, in \si{\per\centi\metre}. In one second the light
travels a distance $c$, and so passes $c\tilde\nu$ crests, each one cycle
of the oscillation, so its frequency is $\nu = c\tilde\nu$. In the units of this book, $c =
\SI{2.99792458e8}{\metre\per\second} = \SI{2.99792458e10}{\centi\metre\per
\second} = \SI{2.99792458e-5}{\centi\metre\per\femto\second}$, the last
step dividing by $10^{15}$ femtoseconds per second, so
\begin{equation*}
   \nu = 2.99792458\times10^{-5}\,\tilde\nu ,
   \qquad
   \text{period} = \frac1\nu , \qquad \omega = 2\pi\nu ,
\end{equation*}
with $\tilde\nu$ in \si{\per\centi\metre}, $\nu$ in cycles per
femtosecond and the period in femtoseconds. In \texttt{mdlab} this speed is
\texttt{units.C\_CM\_PER\_FS}, computed from the exact SI value.

\subsection*{The fastest motions}

The vibrational frequencies of small molecules have been measured
carefully and collected in tables. \Cref{tab:os-vibrations} lists a few
from Shimanouchi's compilation\cite{Shimanouchi1972}, as given by the NIST
Chemistry WebBook\cite{NISTWebBook}, with their periods. A water molecule's
O--H stretches repeat every \SI{9}{\femto\second}, and the C--H stretches
of methane every \SI{11}{\femto\second}; the bend of water takes
\SI{21}{\femto\second}, more than twice as long as its stretches, and the
stretch of the C--C bond in ethane takes \SI{34}{\femto\second}.
Stretches of bonds to hydrogen are the fastest motions in molecules that
contain hydrogen,
for the reason found in \cref{sec:os-bodies}: the hydrogen does almost all
the moving, so the reduced mass is small and $\omega = \sqrt{k/m_{\mathrm r}}$ is
large.

\begin{table}[htb]
\centering
\caption{Measured vibrational wavenumbers of three free (gas-phase)
molecules\cite{Shimanouchi1972,NISTWebBook}, with the frequency and period
they imply.}
\label{tab:os-vibrations}
\begin{tabular}{llccc}
\toprule
molecule & motion & $\tilde\nu$ / \si{\per\centi\metre}
   & $\nu$ / \si{\per\femto\second} & period / \si{\femto\second}\\
\midrule
water & antisymmetric O--H stretch & 3756 & 0.11260 & 8.88\\
water & symmetric O--H stretch & 3657 & 0.10963 & 9.12\\
methane & C--H stretch & 3019 & 0.09051 & 11.05\\
methane & C--H stretch, all in step & 2917 & 0.08745 & 11.44\\
water & bend & 1595 & 0.04782 & 20.91\\
ethane & C--C stretch & 995 & 0.02983 & 33.52\\
\bottomrule
\end{tabular}
\end{table}

No comparably simple measured value is collected here for a lithium atom
in graphite. The model surface of \cref{sec:os-small} gives a period of
\SI{170.3}{\femto\second}, a wavenumber of \SI{196}{\per\centi\metre}, but
it was built for illustration; Chapters~16 and 17 measure the motion of
lithium in graphite with a fitted model of the forces.

\subsection*{The step of a simulation}

A computer that follows atoms advances their positions and velocities by
repeated short steps, as in \cref{sec:ne-equations}. \Cref{sec:mo-atoms}
showed that a sinusoid sampled every $\tau$ is described accurately only
when $\omega\tau$ is small, and the same holds for a motion advanced in
steps. A simulation of water must therefore take steps $\tau$ with
$\omega\tau$ small for its fastest vibration, of period
\SI{8.88}{\femto\second} and so $\omega = 2\pi/8.88 =
\SI{0.708}{\radian\per\femto\second}$: steps short compared with $1/\omega
= \SI{1.41}{\femto\second}$, even if every other motion in it is a hundred
times slower. Chapter~9 works out how much
shorter, and Chapter~11 shows how the fastest vibrations can be frozen out
so that longer steps become safe. The learned propagators planned for Part~VI take
a different route: trained on such simulations, they predict the state of
the atoms a single long step ahead.

\Needspace{9\baselineskip}
\begin{bridge}
The vibrational frequencies computed by electronic-structure codes come
from the eigenvalue problem of \cref{sec:os-modes}: the Hessian of
the Born--Oppenheimer energy, usually by finite differences of the forces,
is mass-weighted and diagonalised. The harmonic frequencies so obtained
neglect the anharmonicity of \cref{sec:os-anharmonic}, which is one
reason they differ from the measured fundamentals of
\cref{tab:os-vibrations}.
\end{bridge}

\begin{selfcheck}
A vibration has a wavenumber of \SI{1000}{\per\centi\metre}. What is its
period in femtoseconds?
\end{selfcheck}
